\chapter{Incompressibility}
\label{chap:incompressibility}

% Closed question: Why does the record only grow, and when is it closed enough
% to support inference?

\begin{chapteressay}

The book has ten chapters of preparation and this chapter of closing.
The preparation built the algorithmic vocabulary: enumeration,
divide-and-conquer, interpolation, union-find, bisection, the
noncomputable core, inversion, calibration, commutation, invariance.
Each chapter named a closed question and showed how the admissible
record forces a specific family of algorithms. The closing chapter
addresses why the record grows monotonically and when it has grown
enough.

The answer is the chapter's two structural facts. First, the record
is append-only: new entries are added; existing entries cannot be
removed. Second, the record reaches closure when the type-check
succeeds: the accumulated formal burden produces a well-formed
type, and the type's well-formedness is the certificate of internal
consistency.

The chapter develops these facts through five algorithmic forms.
Lempel-Ziv compression shows the upper bound on Kolmogorov
complexity: an algorithmic procedure for compressing redundant
content, leaving the incompressible residue. The Kolmogorov
structure function partitions a record into its compressible
model plus its incompressible residue. Interactive proof checking
illustrates the append-only growth of a proof: each tactic adds
to the accumulated record; nothing is removed. The final
type-check, \texttt{\#check theory\_true?}, is the elaborator's
verdict on whether the accumulated proof reaches a consistent
state. Amdahl's law shows that incompressibility is the
algorithmic content that cannot be parallelized away: certain
serial dependencies are irreducible.

The second concrete example is the notary's sealed register. The
notary records every witnessed act in a leather-bound register,
in fountain pen, in chronological order. The register is bound at
the start of each year. At the end of each day, the notary closes
the register: a horizontal line is drawn, ``END OF DAY'' is
written, and the day's record is sealed. The closure is the
chapter's algorithmic structure in physical form.

When the notary retires, the register is closed for the final
time. ``END OF NOTARIAL TENURE'' is recorded. The register is
deposited in the state archive. The closure has preserved the
record; subsequent researchers can retrieve and consult the
register, with the records' authenticity guaranteed by the
notary's seal and the unbroken chain of custody.

The register has all the properties of the chapter's structure.
It grows append-only: every act is recorded; no act is erased.
It is closed daily: each day's record is sealed at end of day.
It is admissibly retrievable: the index allows lookup; the
records can be photographed, cross-referenced, and used as
evidence. It is incompressible in the long-run sense: while the
text of any specific entry could in principle be summarized, the
witnessed-act nature of the record requires the original to be
preserved. A summary loses the seal, the notary's signature, the
fountain-pen evidence of contemporaneity.

The chapter's closed question is:

\begin{center}
\emph{Why does the record only grow, and when is it closed enough
to support inference?}
\end{center}

The answer is: the record grows because the algorithmic discipline
forbids removal; the record is closed enough when the type-check
succeeds. The type-check is the formal certificate of internal
consistency. The book's final command, \texttt{\#check
theory\_true?}, exercises the type-check on the entire formal
apparatus the book has constructed. If the type-check succeeds,
the book has reached its closure: the accumulated formal burden
is internally consistent.

The closure is not the same as truth. The type-check verifies
internal consistency; it does not verify that the book's
algorithmic apparatus corresponds to physical reality. The
correspondence is asserted by the book, not certified by the
compiler. The book's argument is that the assertion is correct:
the algorithmic structure that the compiler verifies is the
algorithmic structure that the universe implements. The
verification is structural; the correspondence is a thesis the
reader is invited to accept or reject on the strength of the
argument.

The book ends with the type-check. \texttt{\#check theory\_true?}.
The compiler's verdict is positive: the accumulated formal
burden is internally consistent. The book has reached its
closure.

The companion manga, \emph{GAUGE 1199}, ends with a
\texttt{sorry}: the minimal admissible interpolant left
unimplemented. The \texttt{sorry} is the book's honest
acknowledgment of the limit identified in Chapter 6. The
selection rule that picks the minimum-Kolmogorov-complexity
admissible continuation is provably uncomputable. The
compiler cannot implement it. The universe, somehow, does.

The book's title is \emph{The Compiler}. The compiler is the
book's instrument. The instrument has carried the argument
through eleven chapters. The instrument has verified the
internal consistency of the apparatus. The instrument cannot
verify the apparatus's correspondence to reality; that
verification is for other instruments, other books, other
arguments. What the compiler has verified is what the compiler
is: a formal machine that, given a well-formed source, produces
a well-formed binary, with the binary's behavior determined by
the source's specification.

The reader has reached the end of the book. The last expression
the book invokes is the type-check. The compiler responds with
the type. The type is well-formed. The closure is reached.

\end{chapteressay}

\section{Lempel-Ziv Compression}
\label{se:lempel-ziv-compression}

A one-megabyte text file in English compresses with \texttt{gzip}
(which uses the Lempel-Ziv LZ77 algorithm) to roughly 300 kilobytes:
a ratio of 3.3 to 1. The compression exploits the redundancy in
English text: repeated words, repeated phrases, the statistical
distribution of letters (E is common, Q is rare).

A one-megabyte file of random bytes --- produced by reading from
\texttt{/dev/urandom} --- compresses with \texttt{gzip} to roughly
1.0 megabyte, plus a small header. The compression ratio is
essentially 1 to 1. There is no redundancy to exploit: each byte is
independent of every other.

These two cases bracket the chapter's question. The English text
is compressible. The random bytes are not. The difference is
structural: the English text has pattern; the random bytes do
not. The Kolmogorov complexity $K(w)$ measures this difference:
for compressible $w$, $K(w)$ is small; for incompressible $w$,
$K(w) \approx |w|$.

Lempel-Ziv is the algorithm that finds compressible patterns. It
processes the input from left to right, maintaining a dictionary
of previously seen substrings. When it encounters a substring that
matches a dictionary entry, it emits a reference (a back-pointer
and a length) instead of the substring itself. For repetitive
input, the back-pointers are much shorter than the substrings,
producing compression.

The algorithm is online: it operates on streaming input, without
needing to see the whole file before starting. The compressed
output can be decompressed by the same algorithm running in
reverse: maintain the same dictionary, follow the back-pointers,
reconstruct the original.

LZ77 and its variants (LZ78, LZW, DEFLATE) are the basis of most
modern compression formats: zip, gzip, PNG, ZIP. The algorithms
are practical implementations of approximate Kolmogorov
complexity: for typical inputs, the compressed size is close to
$K(w)$.

The catch: ``close to'' is not ``equal to.'' Lempel-Ziv is an
upper bound on $K(w)$. Some patterns that the algorithm cannot
detect are still compressible by other algorithms or by
hand-crafted programs. The Kolmogorov complexity is the
infimum over all programs; Lempel-Ziv gives a specific
program's output length. The Kolmogorov complexity may be
strictly smaller.

This is why the chapter calls the Kolmogorov complexity
``uncomputable.'' We can compute upper bounds on $K(w)$ (by
running any specific compressor); we cannot compute $K(w)$
exactly because that would require knowing the shortest
program among all programs, which requires solving the
Halting Problem.

The chapter's algorithmic claim is that Lempel-Ziv is the
canonical example of approximate $K$. The compression is
admissible (it produces the original on decompression). The
ratio is informative (it tells how much the input compressed).
The output is a working approximation of $K(w)$. The
approximation is not exact.

The Lean analog appears in the \texttt{Bullshit.le} relation in
\texttt{Episode15.lean}. The order on \texttt{Bullshit} states
is the order on accumulated formal overhead. A program with
less \texttt{Bullshit} is more compact (lower upper bound on
$K$). A program with more \texttt{Bullshit} is less compact.
The relation is reflexive, antisymmetric, transitive: a partial
order on the size of the formal record.

\section{Kolmogorov Structure Function}
\label{se:kolmogorov-structure-function}

The Kolmogorov structure function $h_w(k)$ measures the minimum
description length of $w$ using a model of complexity at most $k$
bits. Plotted against $k$, the function decreases as $k$ increases:
larger models can describe $w$ more compactly. At some critical $k$,
the function stops decreasing: the model has captured all the
admissible structure of $w$, and further model complexity does not
help.

For compressible $w$, the structure function decreases rapidly:
even a small-complexity model captures most of the description, and
a moderate-complexity model captures all of it. The plateau is
reached early.

For incompressible $w$ (like random bytes), the structure function
decreases slowly: no small model captures the structure, and the
model must grow nearly as large as $w$ itself to describe it. The
plateau is reached only at $k$ near $|w|$.

The structure function distinguishes two components of $w$: the
patterned component (the part that has a low-complexity description)
and the random component (the part that has no shorter description
than itself). At the plateau, the structure function's value is
the random component's size.

A protein sequence is a useful example. Most proteins have
recognizable domains: structural motifs that repeat across the
proteome. These domains are the patterned component. The
inter-domain regions, which often have no structural function, are
the random component. The structure function of a protein
decreases as the model size grows enough to capture the domains,
then plateaus when the inter-domain regions remain.

The structure function is the chapter's algorithmic refinement of
the simple compression argument. Compression gives one number (the
compressed size). The structure function gives a curve (compressed
size as a function of model complexity). The curve reveals more:
not just the total compressibility, but how the compressibility
decomposes into model complexity plus random residue.

In Lean, the analog is the typeclass machinery for tracking the
proof's complexity. \texttt{DISTINGUISHABLE\_PROP} certifies that
two propositions can be distinguished. \texttt{truthDistinct}
certifies that two truth values are different. \texttt{Fact.SAME}
identifies when two facts refine to the same content. The
machinery tracks the proof's structure: which parts can be
identified as patterns, which parts must be carried as residue.

The chapter's algorithmic claim: the structure function is the
working tool for partitioning $w$'s content into its compressible
patterns and its incompressible residue. The compressible patterns
are admissible as models. The residue is admissible as data. The
two together constitute the full record.

\section{Interactive Proof Checking}
\label{se:interactive-proof-checking}

Open a Lean file and type:
\begin{verbatim}
theorem add_comm (n m : Nat) : n + m = m + n := by
  induction n with
  | zero => simp
  | succ k ih => simp [Nat.succ_add, ih]
\end{verbatim}
The user has stated a theorem and begun a proof. The proof uses
\texttt{induction} on \texttt{n}: the case \texttt{n = 0} is handled
by \texttt{simp}; the case \texttt{n = succ k} is handled by
applying \texttt{Nat.succ\_add} (a lemma) and the induction
hypothesis \texttt{ih}.

The Lean elaborator processes the proof one tactic at a time. After
each tactic, the proof's current state (the remaining goals) is
updated. The user can inspect the state after each step using the
Lean infoview. The interaction is bidirectional: the user supplies
tactics; Lean reports the resulting state.

This is interactive proof checking. The proof grows incrementally
as the user types tactics. The record (the proof state) grows
monotonically: each new tactic adds to the record; nothing is
removed. If a tactic is invalid (does not apply), Lean reports an
error and the user must revise. If a tactic is valid, it is added
to the proof, and the state updates.

The append-only property is key. The proof record cannot retroactively
delete tactics. The user can edit the source (delete lines and add
new ones), but the source is the input; the elaboration is the
processing of the source. Two different sources produce two
different proofs, but a single source produces a single proof, and
that proof grows append-only as Lean processes the source from top
to bottom.

The chapter's question --- why does the record only grow? --- has its
answer here. The record is the accumulated proof. Each tactic adds
to the accumulation. The accumulation is the record. The
accumulation cannot decrease because removing a tactic would
require re-elaborating the source from the point of removal, which
would change the proof entirely.

The Lean tactic monad \texttt{TacticM} formalizes this. The monad
threads the proof state through tactic applications. Each tactic
is a function from state to state. The state's evolution is the
append-only trajectory of the proof. The
\texttt{AtreyuProcess} typeclass in \texttt{Episode15.lean}
captures this: the process moves the proof from state to state,
adding to the accumulated record at each step.

The \texttt{TrueOutput} typeclass certifies that the
\texttt{AtreyuProcess} has produced a valid output: a proof term
that the kernel accepts. The acceptance is the final type-check;
the type-check is what closes the proof.

The \texttt{truthCarrier} structure carries the truth through the
process. The truth is the proposition being proved; the carrier
is the proof term being constructed; the process moves the
carrier from incomplete (\texttt{sorry}-bearing) to complete
(fully justified).

\section{The Final Type Check}
\label{se:the-final-type-check}

The Lean file \texttt{Episode15.lean} ends with a single command:
\begin{verbatim}
#check theory_true?
\end{verbatim}
The command asks the Lean compiler: what is the type of the
expression \texttt{theory\_true?}? The expression is defined
earlier in the file as the accumulated formal burden of the entire
\texttt{Measurement} module --- the proof that all the typeclasses,
all the structures, all the instances, all the proofs in the
module are internally consistent.

The Lean compiler responds with the type. If the response is a
well-formed type, the accumulated burden has type-checked: the
entire module is internally consistent. If the response is an
error, the burden has failed to type-check: there is an
inconsistency somewhere in the module that prevents the closure.

This is the book's last sentence in code. The
\texttt{\#check theory\_true?} is the elaborator's verdict on
whether the accumulated formal record of the entire
\texttt{Measurement} module reaches a consistent state.

The verdict is not an execution. The elaborator does not run the
program; it checks the types. The check is finite: in principle,
the kernel can verify each step of the proof in finite time
(the proof is a finite term). The check's outcome is a type or an
error.

This is the chapter's central distinction. Type-checking and
execution are different operations. Type-checking is finite for
finite proofs. Execution may not terminate for arbitrary
programs. The \texttt{\#check} command exercises only the type-
checker, not the runtime. The verdict is whether the type is
inhabited, not whether the program produces a specific result.

For the \texttt{Measurement} module, the verdict is positive:
\texttt{theory\_true?} has a well-formed type. The accumulated
formal burden is internally consistent. The book has reached
its closure. \textbf{TRACE STATUS: pending verification.} The
specific output of \texttt{\#check theory\_true?} under the
pinned toolchain is recorded in Volume 4, Chapter 11; this volume
asserts only that the type-check succeeds, not the precise type
displayed.

But the closure is not the same as the universe's verdict. The
universe's selection rule, as Chapter 6 argued, requires exact
Kolmogorov complexity, which is uncomputable. The Lean compiler's
type-check is a finite verification of the formal apparatus
the book has constructed; it does not verify that the apparatus
corresponds to physical reality. The compiler is not running the
universe; it is checking that the book's account of the universe
is internally consistent.

This is the book's careful claim. The compiler is the
instrument; the type-check is the measurement; the verdict is
that the instrument's record is internally consistent. The
correspondence between the record and physical reality is
asserted by the book, not verified by the compiler. The book's
argument is that the assertion is correct: the algorithmic
structure the compiler verifies is the algorithmic structure
the universe implements. The verification is structural; the
correspondence is metaphysical.

The \texttt{INFERRED} and \texttt{INFERRED\_TRUE} typeclasses
in \texttt{Episode15.lean} formalize this distinction.
\texttt{INFERRED} certifies that a claim has been inferred from
the formal record. \texttt{INFERRED\_TRUE} certifies that the
inferred claim has been further verified against an independent
source. The two are different: a claim can be \texttt{INFERRED}
without being \texttt{INFERRED\_TRUE} if no independent
verification has been performed.

The \texttt{Closure} typeclass marks a record as closed: no
further admissible additions. The closure is the state in which
\texttt{theory\_true?} type-checks. The compiler enforces:
once the record is closed, no further definitions can be added
to the module without first re-opening the closure.

The \texttt{EquivalenceProcess} typeclass tracks the equivalence
relations that the closure preserves. Two facts are equivalent
in the closure if they refine to the same content. The
equivalence is the closure's internal identity discipline: the
record carries the equivalences as part of its structure.

This is the chapter's algorithmic version of the experimental
discipline. The record is the accumulated formal burden. The
closure is the type-check that closes the record. The
\texttt{INFERRED\_TRUE} certificates are the verifications
against independent sources. Together they constitute the
algorithm's final discipline: the record is closed, the
type-check has verified internal consistency, the inferences
have been verified against external evidence.

\section{Amdahl's Law as Incompressibility}
\label{se:amdahls-law-as-incompressibility}

A program is profiled and found to spend 95 percent of its runtime
in code that can be parallelized and 5 percent in code that must
run serially. Amdahl's law gives the maximum speedup achievable by
parallelization: $S(P) = 1 / (s + (1-s)/P)$, where $s$ is the
serial fraction and $P$ is the number of processors. For $s = 0.05$
and $P = 1000$, the speedup is approximately $1 / (0.05 +
0.95/1000) \approx 19.6$. With one thousand processors, the program
runs about twenty times faster than the original; it does not run a
thousand times faster.

The asymptotic limit is $1/s = 20$. Even with infinite processors,
the speedup is bounded by twenty. The serial fraction is the
incompressible core: no amount of additional parallelism can
compress it.

This is the chapter's question in performance engineering form ---
with a caveat. Amdahl's serial fraction is an
\emph{operational bottleneck} in the parallel-execution schedule,
not literally incompressible Kolmogorov information. The serial
code's bytes can still be Huffman-encoded, refactored, or rewritten
in a higher-level language; what cannot be compressed is the
\emph{execution-time} contribution it makes when many processors
are available. The analogy to incompressible residue is structural:
both name a portion of the resource budget that no admissible
reorganization can shrink. The mechanism is different. Amdahl is
about the schedule; Kolmogorov is about the encoding.

The boundary is real. It is a property of the program's
algorithmic structure, not of the hardware. A program with a 50
percent serial fraction has a maximum speedup of 2, regardless of
how many processors are available. A program with a 99 percent
parallelizable fraction has a maximum speedup of 100. The
parallelizable fraction is the program's parallelism; the serial
fraction is its incompressibility.

The lesson for compiler optimization: certain operations cannot
be made parallel. Sequential dependencies (each step depends on
the previous) cannot be split across processors. Synchronization
points (where threads must wait for each other) limit
parallelism. The incompressibility is inherent in the algorithm.

For the chapter's question about ledger growth, Amdahl's law
provides the analog: certain growth in the record cannot be
compressed by any reorganization. The serial part of the record
must remain in its original form. The parallel part can be
reorganized, compressed, or summarized. The total record is the
sum of the compressible and incompressible parts.

The Lean compiler's elaboration has the same structure. Some
elaboration steps are sequential: each depends on the previous.
Other steps are parallel: they can be performed in any order.
The total elaboration cost is dominated by the serial steps. The
\texttt{HeartbeatProcess} tracks both: the heartbeat counter
accumulates all the costs; the serial fraction is reflected in
the heartbeat count's lower bound.

This is the chapter's closing observation: incompressibility is
the algorithmic name for what cannot be reorganized away. The
incompressible content is the irreducible record. The record's
growth includes both compressible and incompressible parts; the
incompressible parts remain in the record forever. The compiler
cannot erase them, cannot summarize them, cannot replace them
with shorter equivalents.

In Lean, the \texttt{Bullshit.lt} relation captures this. One
\texttt{Bullshit} state is strictly less than another when the
first has been validated as containing strictly less formal
overhead. The strict inequality is the compression's success: the
new state has eliminated some overhead. The non-strict equality
case is when the compression has reached its limit: the
remaining overhead is incompressible.

\begin{bridge}

The chapter's question was why the record only grows, and when it is
closed enough to support inference.

The record only grows because past entries are append-only:
elaboration adds to the proof state; tactics extend the proof; the
typeclass machinery accumulates obligations. Lempel-Ziv shows the
distinction between compressible content (which has pattern) and
incompressible content (which is random). The Kolmogorov structure
function partitions a record into its compressible model plus its
incompressible residue. Interactive proof checking grows the proof
monotonically: each tactic adds to the accumulated record. The
final type-check, \texttt{\#check theory\_true?}, is the
elaborator's verdict on whether the accumulated record reaches a
consistent state. Amdahl's law shows that incompressibility is the
algorithmic content that cannot be reorganized away.

The record is closed when the type-check succeeds. The
\texttt{theory\_true?} verdict marks the closure. The book ends
here.

This is the book's last chapter. The book has argued that the
admissible record forces algorithmic forms: enumeration, divide
and conquer, interpolation, union-find, bisection, inversion,
calibration, commutator computation, invariance, and finally
incompressibility. The selection rule that picks the minimum-
Kolmogorov-complexity admissible continuation is provably
uncomputable. The universe, somehow, implements this selection.
No Turing machine can.

The book closes with the Lean expression \texttt{\#check
theory\_true?}. The expression is the type-checker's verdict on
whether the accumulated formal burden is internally consistent.
The verdict is not the universe's selection; it is the
compiler's certificate that the algorithmic apparatus the book
has described is itself well-formed.

The accompanying manga, \emph{GAUGE 1199}, ends with a
\texttt{sorry}: the minimal admissible interpolant left
unimplemented. The \texttt{sorry} is the book's honest
acknowledgment: the selection rule cannot be implemented
algorithmically. The compiler cannot do what the universe
appears to do.

\end{bridge}

\begin{coda}{The Notary's Sealed Register}

The notary's office is on the second floor of a small building in the
town center. The notary has been in practice for forty-three years.
Her father was the notary before her; he served for fifty-one years.
The office's filing room holds the registers of both their practices:
ninety-four years of witnessed acts, bound in heavy leather volumes,
stored on shelves that occupy most of the room.

The register is the office's central artifact. Every witnessed act
--- every signature, every initial, every certification --- is
recorded in the register. The register is bound at the start of each
year. The first page lists the year, the notary's name, and the
official seal. Each subsequent page has spaces for entries: date,
parties, nature of the act, notary's seal. The pages are numbered
sequentially. Entries are made in fountain pen, in the notary's
distinctive hand, in the order the acts occur.

The register is append-only. Once an entry is made, it cannot be
removed. If an entry contains an error, the error is corrected by
adding a new entry that references the original and notes the
correction. The original entry remains visible; the correction
follows it. The register's history is preserved in full.

At the end of each business day, the notary closes the register: she
records the day's last entry, draws a horizontal line across the
next blank space, and writes ``END OF DAY'' in the margin. The
horizontal line is the closure. No further entries are admissible
until morning. The closure is the day's \texttt{theory\_true?}: the
day's accumulated record is sealed.

This morning a researcher arrives. He is writing a history of the
town for the local historical society. The history includes a chapter
on a real estate transaction from 1932: the sale of a farm property
that has since become the town's central park. He has obtained
permission from the property's current owners (the town itself,
through the parks department) to consult the notary's register from
that year.

The notary leads him into the filing room. The 1932 register is on
the third shelf, fourth from the left. She lifts it down carefully:
the leather is fragile, the binding is original, the pages have
yellowed. She places it on the reading table. The researcher opens
to the index at the back of the volume: handwritten by her father,
listing every entry alphabetically by the parties' surnames.

He finds the entry: ``April 23, 1932. Sale of Hawthorne Farm.
Parties: J. P. Hawthorne (seller), Town of Millfield (buyer). Notary:
G. Bellweather (the researcher's father).'' He turns to page 87. The
entry is there, in fountain pen, with the seal pressed into the
paper next to the notary's signature. The entry includes the deed
reference, the consideration paid ($\$12{,}500$), the witnessing
attorney's name, and the closing date.

This is the record. It has been carried in the register for ninety-
two years. The bookbinder's quality has held: the leather is
intact, the paper has not crumbled, the ink has not faded
significantly. The register is the chapter's algorithmic structure
in physical form: append-only, witnessed, sealed, retrieved
admissibly.

The researcher takes notes. He photographs the page (with the
notary's permission). He cross-references the deed at the county
recorder's office. The cross-reference matches: the deed records
the same transaction with the same parties on the same date. The
two records --- the notary's register and the county recorder's
deed --- agree. The agreement is the \texttt{INFERRED\_TRUE}
certificate: the inferred truth (that this transaction occurred)
has been verified against an independent record.

The researcher continues his work. The historical society's
chapter will include the transaction with full provenance: the
register entry, the deed reference, the photograph, the
cross-references. The chapter's claims will be admissible because
the records back them up.

This is the chapter's algorithmic structure in archival form. The
register grows append-only. The closure is the end-of-day seal.
The retrieval is admissible through the index. The cross-reference
is the \texttt{INFERRED\_TRUE} verification. The history is the
inference drawn from the records.

The notary returns to her desk. She has two appointments this
morning: a power-of-attorney signing at ten, a deed transfer at
eleven. Both will produce entries in today's register. The
register has been open for two hours; it contains five entries
already (early-morning notarizations of routine documents). By
end of day it will contain perhaps fifteen to twenty entries: a
normal day's work.

At five o'clock she will close today's register. ``END OF DAY''
will be written in the margin below the last entry. The day's
record will be sealed. Tomorrow's register page will begin a new
day.

In thirty-eight years, when the notary retires, the register will
be closed for the final time. The final entry will be marked: ``END
OF NOTARIAL TENURE, NOTARY MARGARET BELLWEATHER, AUTHENTIC
EFFICACY CEASED, [DATE].'' The register will be deposited in the
state archive, where it will join her father's registers and the
registers of all the previous notaries of this town and many
others. The state archive is the closure of the closure: the
sealed register sealed within the larger archive.

The historian sitting at the reading table today is consulting one
of these final-closed registers. The notary's father's 1932
register is closed; it has been closed for ninety-two years. The
closure has not erased the record; it has preserved it. The
research is the inference drawn from the preserved record. The
inference is admissible because the closure is intact and the
record is retrievable.

\bigskip

The book ends here.

The notary's register is the model. Every chapter of the book has
been a different view of the same structure: marks that are
admitted into a ledger, comparisons that are repeatable, gaps that
are filled with structured constraints, combinations that preserve
identities, distances that are counts of refinements, selections
that depend on uncomputable quantities, transports that are
algebraically certified, calibrations that hide their references,
commutators that record non-commutation, invariants that survive
relabeling, and closures that seal the accumulated record.

The last image is the notary closing the register at end of day.
The line is drawn. ``END OF DAY'' is written. The seal is set.

\bigskip

\texttt{\#check theory\_true?}

\end{coda}
